\documentclass[11pt]{article}

% ---- standalone supplement; compiles on its own (e.g. in Overleaf) ----
% It does not \input or depend on the main paper sources.

\usepackage[margin=1in]{geometry}
\usepackage{amsmath,amssymb}
\usepackage{booktabs}
\usepackage[table]{xcolor}
\usepackage{microtype}
\usepackage[colorlinks=true,allcolors=blue!55!black]{hyperref}

\newcommand{\R}{\mathbb{R}}
\newcommand{\Nn}{\mathcal{N}}
\DeclareMathOperator{\clip}{clip}

\title{\textbf{Supplement: A Motivation--to--Color Readout for the\\
Participation Diffusion Model}}
\author{Hudson Dong \thanks{Class of 2027, Stuyvesant High School, 345 Chambers St, New York, NY 10282. Email: \texttt{hdong70@stuy.edu}}}
\date{}

\begin{document}
\maketitle
\vspace{-3em}

\begin{abstract}
\noindent
This short note formalizes the visual encoding used by the interactive companion to
\emph{Participation Diffusion Dynamics in Influence-Driven Networks}. It defines the map from a
student's scalar motivation state to a display color, records the perceptual rationale for the
choice of color scale, distinguishes core-network nodes from public-body display dots, and defines
the aggregate and ensemble readouts that summarize the motivational state of the network at a
glance. The note is self-contained and introduces no changes to the underlying model.
\end{abstract}

\section{Setup}

Let the club consist of $n$ students indexed by $i \in \{1,\dots,n\}$, embedded in a weighted
directed influence network with in-neighborhoods $\Nn(i)$. At each discrete time (week) $t$, the
model assigns student $i$ a bounded \emph{motivation} state
\begin{equation}
  m_i(t) \in [0,1],
  \label{eq:m}
\end{equation}
together with a \emph{loyalty} stock $L_i(t)\in[0,1]$ and a \emph{fatigue} stock $F_i(t)\in[0,1]$.
We treat \eqref{eq:m} as given by the model dynamics; the present note concerns only how
$m_i(t)$ and $L_i(t)$ are rendered, and how they are summarized across $i$.

Because $m_i(t)$ is already normalized to the unit interval, no instance-dependent rescaling is
required: the encoding below is a fixed function of $m_i(t)$, so colors are comparable across
students, across weeks, and across separate simulation runs.

\section{The color map}

\paragraph{Anchor stops.}
Let $\mathsf{c}_0,\dots,\mathsf{c}_K \in [0,255]^3$ be $K{+}1$ anchor colors in RGB space placed at
equally spaced levels $\ell_k = k/K$. We use the $K{=}5$ \emph{plasma} anchors of
Table~\ref{tab:stops}, a perceptually ordered sequence running from dark blue through magenta and
orange to bright yellow.

\begin{table}[h]
\centering
\renewcommand{\arraystretch}{1.15}
\begin{tabular}{cccccc}
\toprule
$k$ & level $\ell_k$ & $R$ & $G$ & $B$ & swatch \\
\midrule
$0$ & $0.0$ & $13$  & $8$   & $135$ & \cellcolor[RGB]{13,8,135}\hspace{2.4em} \\
$1$ & $0.2$ & $126$ & $3$   & $168$ & \cellcolor[RGB]{126,3,168}\hspace{2.4em} \\
$2$ & $0.4$ & $177$ & $42$  & $144$ & \cellcolor[RGB]{177,42,144}\hspace{2.4em} \\
$3$ & $0.6$ & $225$ & $100$ & $98$  & \cellcolor[RGB]{225,100,98}\hspace{2.4em} \\
$4$ & $0.8$ & $252$ & $166$ & $54$  & \cellcolor[RGB]{252,166,54}\hspace{2.4em} \\
$5$ & $1.0$ & $240$ & $249$ & $33$  & \cellcolor[RGB]{240,249,33}\hspace{2.4em} \\
\bottomrule
\end{tabular}
\caption{Plasma anchor stops $\mathsf{c}_k$ at levels $\ell_k = k/5$.}
\label{tab:stops}
\end{table}

\paragraph{Piecewise-linear interpolation.}
Define the color map $\mathcal{C}:[0,1]\to[0,255]^3$ by linear interpolation between adjacent anchors.
For $m\in[0,1]$, let $u = Km$, $k=\lfloor u\rfloor$ (capped at $K{-}1$), and $f = u-k$; then
\begin{equation}
  \mathcal{C}(m) \;=\; (1-f)\,\mathsf{c}_k \;+\; f\,\mathsf{c}_{k+1}.
  \label{eq:cmap}
\end{equation}
A core-network student with motivation $m_i(t)$ is drawn as a disk filled with
$\mathcal{C}\!\big(m_i(t)\big)$. The map $\mathcal{C}$ is continuous and applied componentwise,
so small changes in motivation produce small, smooth changes in color.

\paragraph{Perceptual rationale.}
The plasma sequence is approximately \emph{perceptually uniform} and increases monotonically in
luminance from $\mathsf{c}_0$ to $\mathsf{c}_5$. Two consequences matter for reading the display.
First, ordering is preserved: if $m_i < m_j$ then student $i$ appears uniformly ``cooler/darker''
than student $j$, so relative motivation can be judged by brightness alone. Second, equal
differences in $m$ correspond to roughly equal perceived color differences, so the spatial
\emph{gradient} of color across the network is an honest proxy for the spatial gradient of
motivation, rather than exaggerating or compressing particular ranges. The map degrades gracefully
to grayscale luminance, preserving the ordering for color-vision deficiencies.

\paragraph{Core-network ring and loyalty overlay.}
Every simulated core-network student is marked by an outer ring so that the viewer can distinguish
modeled agents from the wider public-body dots. Motivation remains the fill channel: dormant or
low-motivation students are dark, while motivated students move along the plasma scale. Loyalty is
rendered on a separate, non-competing channel: an amber arc of angular extent $2\pi\,L_i(t)$ is drawn
around the node. Motivation (fill hue and brightness), core membership (outer ring), and loyalty
(arc sweep) are thus separable visual variables. Inactive students --- those whose participation
indicator has lapsed --- are desaturated and dashed, keeping them visible without treating them as
current participants.


\section{Core-network and public-body display layers}

The mathematical model is defined on the simulated core network $V=\{1,\dots,n\}$. The interactive
display may also show a surrounding field of public-body dots representing students outside that core
network. These dots are a presentation layer: they help communicate that the club is embedded in a
larger school population, but they are not additional model agents unless explicitly included in $V$.
For this reason public-body dots are rendered without the core-network ring. If a public-body dot is
shown as ``drawn in,'' that status is a qualitative display readout derived from turnout, events, and
acquaintance reach; it should not be confused with the state variables $m_i,L_i,F_i$ governed by the
model equations.

The display therefore uses a simple visual grammar:
\begin{itemize}
  \item filled disk color $\mathcal{C}(m_i)$: current motivation of a simulated core-network agent;
  \item outer ring: membership in the simulated core network;
  \item amber arc: loyalty stock $L_i$;
  \item dashed/desaturated node: lapsed or inactive core-network agent;
  \item plain faint dot: surrounding public-body display dot.
\end{itemize}

\section{Aggregate reach scalars}

To summarize the network's motivational state at week $t$ in a single number we use the mean
motivation, which we call the \emph{reach} (or affectedness) of the diffusion process:
\begin{equation}
  \bar{m}(t) \;=\; \frac{1}{n}\sum_{i=1}^{n} m_i(t) \;\in[0,1].
  \label{eq:reach}
\end{equation}
Equation~\eqref{eq:reach} is the spatial average of the field that $\mathcal{C}$ colors; it rises as
motivation spreads through $\Nn(\cdot)$ and falls as it drains, and it is directly comparable across
runs because each $m_i\in[0,1]$.

It is sometimes useful to count only students whose motivation has crossed a salience threshold
$\tau\in[0,1]$. Define the \emph{active reach fraction}
\begin{equation}
  \rho_\tau(t) \;=\; \frac{1}{n}\,\bigl|\{\, i : m_i(t) > \tau \,\}\bigr| \;\in[0,1],
  \label{eq:rho}
\end{equation}
the share of the network rendered at or above the color level $\mathcal{C}(\tau)$. The companion uses
a fixed $\tau$ to report a ``reach'' percentage; \eqref{eq:rho} is monotone nonincreasing in $\tau$,
with $\rho_0(t)$ counting all students with any motivation and $\rho_\tau(t)\le \bar m(t)/\tau$ by
Markov's inequality.

A complementary scalar is the cross-sectional spread of motivation,
\begin{equation}
  s(t) \;=\; \sqrt{\tfrac{1}{n}\textstyle\sum_{i=1}^{n}\bigl(m_i(t)-\bar m(t)\bigr)^2},
  \label{eq:spread}
\end{equation}
which is small when the network is uniform (all nodes a single color --- either a dim, collapsing
club or a fully saturated one) and largest during the transient phases when a motivated core
coexists with an unconvinced periphery. Reading $\bar m(t)$ together with $s(t)$ distinguishes a
genuinely thriving, differentiated network from one that is uniformly saturated, even though both can
share a high attendance count.


\section{Sample and ensemble readouts}

The companion provides two numerical views of the same equations. A sample run displays one
realization of the randomized graph and heterogeneous initial conditions. An ensemble run averages
$M$ independent realizations under the same intervention path. For any scalar network quantity
$Z(t)$, such as attendance $A(t)$, mean motivation $\bar m(t)$, or active reach $\rho_\tau(t)$, the
reported ensemble mean is
\begin{equation}
  \widehat Z_M(t)=\frac{1}{M}\sum_{r=1}^{M} Z^{(r)}(t),
  \,\qquad t=0,\dots,T.
  \label{eq:ensembleZ}
\end{equation}
The time scrubber in ensemble mode moves through these averaged weekly states. Uncertainty bands are
empirical quantiles of $\{Z^{(r)}(t)\}_{r=1}^{M}$ and are visual summaries of random-network
sensitivity, not additional dynamical variables.

When the network panel displays averaged nodes, the corresponding value for a core-network index is
\begin{equation}
  \widehat m_i(t)=\frac{1}{M}\sum_{r=1}^{M} m_i^{(r)}(t),
  \label{eq:ensemblem}
\end{equation}
and analogously for $L_i,F_i$, and attendance probability. Such node-level values should be read as
averaged positions in the simulated core network, not as predictions about named students. This
preserves the original model while allowing the interface to show both a mechanism-level sample and a
robustness-oriented ensemble summary.

\section{Summary}

The display encodes the model's per-student motivation $m_i(t)$ through the fixed, perceptually
ordered map $\mathcal{C}$ of \eqref{eq:cmap}, with core-network membership marked by an outer ring
and loyalty $L_i(t)$ on an orthogonal arc channel. The scalars $\bar m(t)$, $\rho_\tau(t)$, and
$s(t)$ of \eqref{eq:reach}--\eqref{eq:spread} summarize, respectively, how far motivation has
spread, what fraction of the network it has reached, and how unevenly it is distributed. The
ensemble quantities \eqref{eq:ensembleZ}--\eqref{eq:ensemblem} average these same model outputs over
randomized graph realizations; they do not add or modify behavioral equations.

\end{document}
